% ============================================================================
% Fichier  : partie1/P01_C01_donnees_traitements.tex
% Titre    : Données, Traitements, Acteurs
% Auteur   : MINKA MI NGUIDJOÏ Thierry Emmanuel
% Version  : 1.0 -- novembre 2026
% Corrigés : corriges/Corr_C01.tex
% ============================================================================
% Sources : loi n° 2024/017, art. 1 à 5, 20, 31, 48, 72 (texte officiel) ;
% ISO/IEC 29100 (cadre de protection de la vie privée, paraphrasé) ;
% Solove (2008) ; Nissenbaum (2010).
% Faits nouveaux (bible) : inventaire des familles de données de NYANGA,
% 7 juillet 2026.
% ============================================================================

\chapter{Données, Traitements, Acteurs}
\label{chap:donnees}
\chaptermark{Données, traitements, acteurs}

\epigraph{%
    « On ne protège bien que ce que l'on sait nommer. »%
}{Principe du cours}

\niveauChapitre{Fondamental}{%
    Lecture du Chapitre~\ref{chap:prologue}
}

\begin{boxobjectifs}
\begin{itemize}
    \item Reconnaître une donnée à caractère personnel, directe ou
          indirecte, selon l'article~5 de la loi n°~2024/017.
    \item Distinguer anonymisation et pseudonymisation, et dire ce que
          chacune change pour l'application de la loi.
    \item Identifier les données sensibles, dont la liste camerounaise est
          plus large que la liste européenne.
    \item Qualifier les acteurs d'un traitement~: personne concernée,
          responsable, sous-traitant, responsables conjoints, destinataire.
    \item Délimiter le champ d'application de la loi et ses exclusions.
\end{itemize}
\end{boxobjectifs}

\bigskip

% ============================================================================
\section*{Avant de commencer : « nos données personnelles, ce sont les clients »}
% ============================================================================

Bafoussam, mardi 7~juillet~2026. Pour sa deuxième journée de terrain,
Mme~BELLO demande à chaque directeur de NYANGA CRÉDIT de lister les données
personnelles que traite son service. Le directeur commercial répond
aussitôt~: « Nos données personnelles, ce sont les clients. » La directrice
des ressources humaines n'a rien noté~; elle pensait que la question ne la
concernait pas.

En fin de journée, l'auditrice a recensé onze familles de données, dont
huit ne figuraient sur aucune liste~: les cautions des emprunteurs, les
prospects démarchés par les agents, les candidats à l'embauche, les
salariés, les images des caméras des agences, les enregistrements du
service client, les journaux de connexion de l'application mobile et les
listes de débiteurs remises au prestataire de recouvrement.

\begin{boxalert}
\textbf{Violation typique.} Réduire les données personnelles aux fichiers
clients, et oublier les salariés, les tiers, les images, les voix et les
journaux techniques.

\smallskip
\textbf{Règle~:} une donnée est personnelle dès qu'elle permet d'identifier
une personne, directement ou indirectement, quel que soit son support et
quelle que soit la personne. Le recensement part des personnes concernées,
pas des applications.
\end{boxalert}

% ============================================================================
\section{La Donnée à Caractère Personnel}
\label{sec:c01-donnee}
% ============================================================================

L'article~5 de la loi définit la donnée à caractère personnel par sa
relation à une personne identifiable, et non par sa nature.

\begin{boxjuris}
\textbf{\loicam{n°~2024/017}, article~5}~\cite{loi2024017}~:
\textit{« Donnée à caractère personnel~: information se rapportant à une
personne permettant de l'identifier directement ou indirectement, notamment
par référence à toute forme d'identifiant ou à un ou plusieurs éléments,
propres à son identité physique, psychologique, génétique, psychique,
culturelle, socio-professionnelle ou économique, notamment un nom, une
photo, une empreinte, une adresse postale, une adresse mail, un numéro de
téléphone, un numéro de sécurité sociale, un matricule interne, un
identifiant numérique, une adresse IP, un identifiant de connexion
informatique, un enregistrement vocal. »}

\smallskip
\textbf{Portée.} La liste des exemples est longue et précise~: une adresse
IP, un matricule interne et un enregistrement vocal sont des données
personnelles par la lettre même de la loi. L'auditeur n'a pas à le
démontrer.
\end{boxjuris}

\subsection{Identifier directement ou indirectement}

Une donnée identifie \termecle{directement} quand elle désigne la personne à
elle seule~: un nom complet, une photographie, une empreinte. Elle identifie
\termecle{indirectement} quand il faut la rapprocher d'autres informations~:
un numéro de compte, rapproché du fichier des comptes~; une adresse IP,
rapprochée du journal de l'opérateur~; la combinaison « femme, 34~ans,
pharmacienne, quartier Tamdja » dans une petite ville. La question n'est pas
de savoir si l'organisation fait ce rapprochement, mais s'il est
raisonnablement possible.

\subsection{Anonymiser n'est pas pseudonymiser}

\begin{boxdef}
\termecle{Anonymisation}~: selon l'article~5, « processus de modification ou
de suppression d'informations à caractère personnel identifiables dans un
ensemble de données, afin de rendre ces données non attribuables à des
individus spécifiques ».

\smallskip
\termecle{Pseudonymisation}~: « technique utilisée pour remplacer des
données d'identification directe par des identifiants artificiels dans le
but de protéger la vie privée des individus » (article~5).
\end{boxdef}

La différence est décisive pour l'auditeur. Des données réellement
anonymisées ne se rapportent plus à une personne identifiable~: elles
sortent du champ de la loi. Des données pseudonymisées restent
personnelles, puisque celui qui détient la table de correspondance peut
rétablir le lien. Un fichier où le nom du client est remplacé par son
numéro de compte n'est pas anonyme~: il est pseudonymisé, et le
numéro de compte est lui-même un identifiant.

\begin{boxinfo}
\textbf{L'anonymisation se prouve.} Une organisation qui déclare « nos
statistiques sont anonymes » doit montrer qu'aucun rapprochement raisonnable
ne permet de retrouver une personne. Supprimer les noms ne suffit pas quand
il reste la date de naissance, le quartier et la profession. Tant que la
preuve manque, l'auditeur cote la proposition « ces données sont anonymes »
\NEK{}, et traite les données comme personnelles.
\end{boxinfo}

% ============================================================================
\section{Les Données Sensibles}
\label{sec:c01-sensibles}
% ============================================================================

Certaines données exposent davantage la personne~: leur divulgation peut
conduire à une discrimination, une persécution ou une atteinte grave à la
réputation. La loi les nomme et les encadre plus strictement.

\begin{boxjuris}
\textbf{\loicam{n°~2024/017}, article~5}~\cite{loi2024017}~:
\textit{« Donnée sensible~: information relative notamment aux opinions et
activités religieuses, philosophiques, politiques, syndicales, aux
transactions bancaires, à l'origine raciale ou ethnique, linguistique ou
régionale, à la vie sexuelle, à la génétique, à la biométrie, à la santé,
aux poursuites judiciaires et aux sanctions pénales. »}

\smallskip
\textbf{Portée.} Le mot « notamment » rend la liste ouverte. Elle comprend
trois catégories que le droit européen ne range pas parmi les données
sensibles~: les transactions bancaires, l'origine linguistique ou régionale,
et les poursuites judiciaires.
\end{boxjuris}

La présence des \textbf{transactions bancaires} dans la liste change tout
pour un établissement financier~: la plupart des données de NYANGA CRÉDIT
et de NOUN PAY est sensible. L'article~48, alinéa~2, en tire la
conséquence~: le traitement des données dans le cadre des transactions
bancaires est interdit « sans autorisation préalable des administrations et
structures compétentes ». L'origine \textbf{linguistique ou régionale}
répond à une préoccupation propre au Cameroun~; un formulaire qui demande la
région d'origine ou la langue maternelle collecte une donnée sensible.

L'article~48, alinéa~1, interdit le traitement des données relatives aux
opinions, à l'origine, à la vie sexuelle, « à la génétique et à la
biométrie en matière de santé ». La lecture de cette fin de phrase n'est
pas tranchée~: la biométrie est-elle interdite seulement en matière de
santé, ou dans tous les cas~? L'auditeur ne tranche pas à la place du juge~;
il signale la difficulté et cote \NEK{} la licéité d'un pointage
biométrique tant qu'aucune autorisation ni aucun texte d'application ne
l'éclaire (Chapitre~\ref{chap:profilage}).

% ============================================================================
\section{Le Traitement et le Fichier}
\label{sec:c01-traitement}
% ============================================================================

\begin{boxdef}
\termecle{Traitement}~: selon l'article~5, « opération ou ensemble
d'opérations effectuées ou non à l'aide de procédés entièrement ou
partiellement automatisés ou non automatisés et appliqués à des données à
caractère personnel, notamment la collecte, l'enregistrement,
l'organisation, la conservation, l'adaptation ou la modification,
l'extraction, la consultation, l'utilisation, la communication par
transmission, diffusion ou toute autre forme de mise à disposition, le
rapprochement, ainsi que le verrouillage, l'effacement ou la destruction ».

\smallskip
\termecle{Fichier}~: « ensemble de données structurées et à caractère
personnel accessibles selon des critères déterminés, que cet ensemble soit
centralisé, décentralisé ou réparti de manière fonctionnelle ou
géographique » (article~5).
\end{boxdef}

Deux conséquences pratiques en découlent. D'abord, le papier compte~: la
loi vise aussi les procédés « non automatisés », et les dossiers de crédit
rangés dans les armoires des agences sont des fichiers. Ensuite, une simple
consultation est un traitement~: l'agent qui ouvre le dossier d'un client
sans raison de service traite des données, et c'est précisément ce que le
CENTRE MÉDICAL MBOLO découvrira dans ses journaux (Chapitre~\ref{chap:securite}).

Un traitement se définit par sa \termecle{finalité}~: l'ensemble des
opérations faites dans un même but. « Gérer les comptes des clients » est
un traitement~; « tenir la base de données Oracle » n'en est pas un, c'est
un moyen technique qui sert plusieurs traitements. Cette distinction
organise tout le recensement du Chapitre~\ref{chap:cartographie}.

\begin{boxobjet}{Gestion des comptes des clients, NYANGA CRÉDIT}
\textbf{Finalité~:} ouvrir, tenir et clôturer les comptes d'épargne et de
crédit. \textbf{Personnes concernées~:} clients, mandataires, cautions.
\textbf{Données~:} identité, pièce d'identité, adresse, téléphone,
profession, revenus déclarés, mouvements des comptes (données sensibles,
article~5). \textbf{Destinataires~:} agences, service crédit, éditeur du
logiciel bancaire, prestataire de recouvrement pour les impayés.
\textbf{Ce que l'audit établit au 7~juillet~:} le traitement existe
(\OUI{})~; son autorisation n'a pas été demandée (\NON{})~; sa durée de
conservation n'est fixée nulle part (\NON{})~; le fondement du traitement
des cautions reste à examiner (\NEK{}).
\end{boxobjet}

% ============================================================================
\section{Les Acteurs du Traitement}
\label{sec:c01-acteurs}
% ============================================================================

\begin{boxdef}
\termecle{Personne concernée}~: « toute personne dont les données à
caractère personnel font l'objet d'un traitement » (article~5).

\smallskip
\termecle{Responsable du traitement}~: « personne physique ou morale qui,
seule ou conjointement, collecte et traite des données à caractère
personnel et en détermine les moyens et les finalités » (article~5).

\smallskip
\termecle{Sous-traitant}~: « toute personne physique ou morale qui traite des
données à caractère personnel pour le compte du responsable du traitement
et sous ses instructions » (article~5).

\smallskip
\termecle{Destinataire}~: « toute personne physique ou morale, habilitée à
recevoir une communication des données à caractère personnel »
(article~5).
\end{boxdef}

La qualification se fait \textbf{opération par opération}, pas acteur par
acteur. L'éditeur du logiciel bancaire de NYANGA CRÉDIT est sous-traitant
quand il héberge les comptes selon les instructions de l'établissement~; il
deviendrait responsable s'il utilisait ces données pour améliorer ses
propres produits sans instruction. Un même acteur peut donc porter deux
qualités pour deux opérations différentes~; le Chapitre~\ref{chap:recherche}
montre que cette règle pose des questions difficiles dans la chaîne
pénale.

\begin{boxjuris}
\textbf{\loicam{n°~2024/017}, article~20}~\cite{loi2024017}~:
\textit{« Le responsable du traitement et le sous-traitant sont assujettis
aux mêmes obligations au titre de l'activité de traitement des données à
caractère personnel. »}

\smallskip
\textbf{Portée.} Le texte camerounais va plus loin que le droit européen,
qui distingue nettement les obligations des deux acteurs. Pour l'auditeur,
le prestataire de recouvrement et l'éditeur du logiciel bancaire répondent
des mêmes obligations que NYANGA CRÉDIT pour les traitements qu'ils
réalisent.
\end{boxjuris}

L'article~31 ajoute les \termecle{responsables conjoints}~: lorsque deux
organisations déterminent ensemble les finalités et les moyens, elles sont
« solidairement responsables » et doivent fixer leurs obligations dans une
convention « rendue publique » (Chapitre~\ref{chap:soustraitants}).

\begin{tableaucours}{L{3.6cm}L{4.4cm}L{4.0cm}}{%
    \tableauHeader{Loi n°~2024/017} &
    \tableauHeader{ISO/IEC 29100} &
    \tableauHeader{RGPD (UE)}
}
Personne concernée & Personne à qui se rapportent les données (\textit{PII
principal}) & Personne concernée \\
Responsable du traitement & \textit{PII controller} & Responsable du
traitement \\
Sous-traitant & \textit{PII processor} & Sous-traitant \\
Destinataire, tiers & Tiers (\textit{third party}) & Destinataire, tiers \\
Autorité de protection & Autorité de contrôle (hors norme) & Autorité de
contrôle \\
\end{tableaucours}
\vspace{-0.8em}
\begin{center}{\small\itshape Tableau~1.1. Les acteurs dans trois vocabulaires}\end{center}

La norme ISO/IEC~29100~\cite{iso29100} fournit un vocabulaire et des
principes communs à toutes les normes de la famille, en particulier
l'ISO/IEC~27701 (Chapitre~\ref{chap:normes}). Le RGPD~\cite{ue2016rgpd}
sert de point de comparaison~: beaucoup d'organisations camerounaises, comme
NOUN PAY, y sont exposées parce que leurs prestataires sont européens.

% ============================================================================
\section{Le Champ d'Application}
\label{sec:c01-champ}
% ============================================================================

L'article~2 soumet à la loi tout traitement effectué par l'État, les
collectivités territoriales décentralisées ou toute personne physique ou
morale~; tout traitement des données de personnes établies, résidant ou en
transit au Cameroun~; tout traitement réalisé par un responsable ou un
sous-traitant établi au Cameroun~; et tout traitement réalisé dans un
territoire où le droit camerounais s'applique. Le critère des personnes
« en transit » est large~: un voyageur dont les données sont traitées
pendant son passage au Cameroun est protégé.

L'article~3 exclut trois catégories~: les traitements d'une personne
physique pour ses activités exclusivement personnelles ou domestiques, à
condition que les données ne soient pas destinées à une communication
systématique, à un tiers ou à la diffusion~; les copies temporaires faites
pour la transmission sur un réseau~; et les traitements à des fins
littéraires, artistiques, archivistiques d'intérêt public, de recherche,
de statistique ou de journalisme, dans le respect des règles
déontologiques de ces professions.

Deux articles renvoient enfin à des textes particuliers~: l'article~4 pour
les traitements des autorités de sécurité et de défense, l'article~72 pour
les traitements relatifs « notamment à la sécurité, à la défense, à la
santé, à la justice et à l'état civil ». La portée de cette réserve est
discutée aux Chapitres~\ref{chap:secteurs} et~\ref{chap:recherche}~: tant
que les textes spécifiques ne sont pas identifiés, l'auditeur applique la
loi générale et signale la réserve.

% ============================================================================
\section{L'Inventaire de NYANGA}
\label{sec:c01-nyanga}
% ============================================================================

\begin{boxscenario}{Mission NYANGA. Onze familles de données, 7 juillet 2026}
\begin{tableaucours}{L{3.6cm}L{4.4cm}cc}{%
    \tableauHeader{Famille} &
    \tableauHeader{Exemples de données} &
    \tableauHeader{Sensible} &
    \tableauHeader{Connue}
}
Clients & Identité, revenus, mouvements & Oui & Oui \\
Cautions & Identité, revenus & Oui & Non \\
Prospects & Nom, téléphone, quartier & Non & Non \\
Débiteurs remis au recouvrement & Identité, montant dû & Oui & Non \\
Salariés & Identité, paie, santé (arrêts) & Oui & Oui \\
Candidats & CV, région d'origine & Oui & Non \\
Images des agences & Visages, heures de passage & Non & Non \\
Appels du service client & Voix, motif d'appel & Non & Non \\
Journaux de l'application & Adresse IP, identifiant & Non & Non \\
Utilisateurs du SI & Comptes, journaux d'accès & Non & Non \\
Fournisseurs (personnes physiques) & Identité, coordonnées bancaires & Oui &
Oui \\
\end{tableaucours}

\medskip
\textbf{Constatations.} Onze familles, dont huit absentes des listes de la
direction. Sept contiennent des données sensibles au sens de l'article~5,
dont quatre à cause des transactions bancaires et une à cause de la
région d'origine demandée aux candidats.

\textbf{Analyse} (confiance élevée). La vision de l'établissement était
organisée par applications et par clients. Les familles oubliées sont
celles dont la finalité n'est pas commerciale~: elles n'ont pas de
« propriétaire » évident dans l'organisation.
\end{boxscenario}

% ============================================================================
\section{Ce que la Notion demande à chaque Métier}
\label{sec:c01-metiers}
% ============================================================================

\begin{boxmetier}{Auditeur de la protection des données}
Recenser à partir des personnes concernées, pas des applications. Pour
chaque famille, demander qui la collecte, pour quoi faire, et où elle va.
Ne jamais accepter « anonyme » sans preuve.
\end{boxmetier}

\begin{boxmetier}{Juriste conformité}
Qualifier chaque acteur pour chaque opération, et relire les contrats à la
lumière de l'article~20~: le prestataire porte les mêmes obligations que
l'établissement.
\end{boxmetier}

\begin{boxmetier}{Responsable de traitement}
Désigner, pour chaque famille de données, un propriétaire dans
l'organisation. Une famille sans propriétaire est une famille que personne
ne protège.
\end{boxmetier}

\begin{boxresume}
\begin{itemize}
    \item Une donnée est personnelle dès qu'elle identifie une personne,
          directement ou indirectement~; la loi cite l'adresse IP, le
          matricule interne et la voix.
    \item Les données pseudonymisées restent personnelles~; seules des
          données dont l'anonymisation est prouvée sortent du champ de la
          loi.
    \item La liste camerounaise des données sensibles inclut les
          transactions bancaires, l'origine linguistique ou régionale et les
          poursuites judiciaires.
    \item Le responsable et le sous-traitant sont assujettis aux mêmes
          obligations (article~20)~; la qualification se fait opération par
          opération.
    \item À NYANGA CRÉDIT, onze familles de données ont été recensées, dont
          huit inconnues de la direction et sept sensibles.
\end{itemize}
\end{boxresume}

\begin{boxinfo}
\textbf{Points de contrôle.} Vous avez compris ce chapitre si vous savez~:
\begin{itemize}
    \item dire pourquoi un fichier pseudonymisé reste soumis à la loi~;
    \item citer trois données sensibles propres à la liste camerounaise~;
    \item qualifier un prestataire de responsable ou de sous-traitant pour
          une opération donnée.
\end{itemize}
\end{boxinfo}

\begin{boxexo}[title={\ding{45}\quad Exercice~1.1. Personnelle, sensible, ni l'une ni l'autre \niveaubase}]
Classez~: (a) le relevé de compte d'un client~; (b) la langue maternelle
d'un candidat~; (c) le chiffre d'affaires global d'une agence~; (d) la
photographie d'un client sur sa carte d'identité~; (e) l'adresse IP d'un
utilisateur de l'application~; (f) le casier judiciaire d'un candidat.
\end{boxexo}

\begin{boxexo}[title={\ding{45}\quad Exercice~1.2. Anonyme ou pseudonyme \niveaumoyen}]
Le service marketing transmet à un cabinet d'études un fichier où chaque
client est désigné par un code, avec son âge, son sexe, son agence, sa
profession et son encours. Le fichier est-il anonyme~? Que faudrait-il
faire pour qu'il le soit, et comment le prouver~?
\end{boxexo}

\begin{boxexo}[title={\ding{45}\quad Exercice~1.3. Qualifier les acteurs \niveaumoyen}]
Pour chaque opération, dites qui est responsable, qui est sous-traitant~:
(a) l'éditeur héberge les comptes et applique les sauvegardes demandées~;
(b) le prestataire de recouvrement appelle les débiteurs selon un script
fourni par NYANGA~; (c) le même prestataire utilise les numéros pour
proposer ses propres services~; (d) un opérateur de téléphonie envoie les
SMS de relance.
\end{boxexo}

\begin{boxexo}[title={\ding{45}\quad Exercice~1.4. Le champ de la loi \niveauexpert}]
\begin{enumerate}[(a)]
    \item Un agent de NYANGA conserve sur son téléphone personnel les
          numéros de ses clients pour les relancer. L'exclusion des
          activités domestiques (article~3) s'applique-t-elle~?
    \item Un journaliste publie des extraits des fichiers exfiltrés en
          mars~2025. Quel régime s'applique à son traitement~?
    \item L'expert judiciaire conserve une copie du dossier \texttt{Clients}
          pour son rapport. Quels articles faut-il lire, et quelle
          difficulté signalez-vous~?
\end{enumerate}
\end{boxexo}

\bigskip

\begin{boxinfo}
\textbf{Pour aller plus loin.}
\begin{itemize}
    \item L'ISO/IEC~29100~\cite{iso29100}, pour le vocabulaire et les onze
          principes de protection de la vie privée.
    \item H.~Nissenbaum, \textit{Privacy in Context}~\cite{nissenbaum2010},
          pour l'idée qu'une donnée est protégée par le contexte dans lequel
          elle circule.
\end{itemize}
\end{boxinfo}

\begin{boxinfo}
\textbf{Transition.} Savoir qu'un traitement existe ne dit pas s'il est
permis. Le Chapitre~\ref{chap:principes} expose les principes que tout
traitement doit respecter, et la règle du consentement, plus stricte au
Cameroun qu'en Europe.
\end{boxinfo}
